\lecture{17}{3. November 2025}{Bending – Deformation of Beams}

\begin{problemwithimage}{f11_43}{11-43}
  A member having the dimensions shown is used to resist an internal bending moment of $M = \qty{90}{\kN.\m}$.
  Determine the maximum stress in the member if the moment is applied (a) about the $z$ axis (as shown), (b) about the $y$ axis.
  Sketch the stress distribution for each case.
\end{problemwithimage}

\begin{derivation}
  We start by calculating the second moments for the rectangular cross section about each of the two axes.
  \begin{align*}
    I_z = \frac{bh^3}{12} &= \frac{\qty{200}{\mm} \cdot \left( \qty{150}{\mm}  \right)^3}{12}  = \qty{56,25e6}{\mm^4}\\
    I_y = \frac{bh^3}{12} &= \frac{\qty{150}{\mm} \cdot \left( \qty{200}{\mm}  \right)^3}{12} = \qty{100e6}{\mm^4}
  .\end{align*}

  \begin{enumerate}[label=(\alph*)]
    \item For the moment about the $z$ axis we use the formula for the maximum stress:
      \begin{equation*}
        \sigma_{\mathrm{max}} = \frac{Mc}{I}
      .\end{equation*}
      For this case, we have $c_z = \frac{h}{2} = \qty{75}{\mm}$ and we get
      \begin{equation*}
        \sigma_{z_{\mathrm{max}}} = \frac{M_z c_z}{I_z} = \frac{\qty{90}{\kN.\m} \cdot \qty{75}{\mm} }{\qty{56,25e6}{\mm^{4}} }  = \qty{120}{\MPa} 
      .\end{equation*}
    \item For the moment about the $y$ axis we use the same formula, however, this time $c_y = \frac{b}{2} = \qty{100}{\mm}$.
      Hence,
      \begin{equation*}
        \sigma_{y_{\mathrm{max}}} = \frac{M_y c_y}{I_y} = \frac{\qty{90}{\kN.\m} \cdot \qty{100}{\mm}}{\qty{100e6}{\mm^{4}} } = \qty{90}{\MPa} 
      .\end{equation*}
  \end{enumerate}
\end{derivation}

\finalresult{
\begin{aligned}
  \sigma_{z_{\mathrm{max}}} &= \qty{120}{\MPa} \\
  \sigma_{y_{\mathrm{max}}} &= \qty{90}{\MPa} 
.\end{aligned}
}

\begin{problemwithimage}{f11_64}{11-63}
  The man has a mass of \qty{78}{\kg} and stands motionless at the end of the diving board.
  If the board has the cross section shown, determine the maximum normal strain developed in the board.
  The modulus of elasticity for the material is $E = \qty{125}{\GPa}$.
  Assume $A$ is a pin and $B$ is a roller. 
\end{problemwithimage}

\begin{derivation}
  The weight of the man is
  \begin{equation*}
    W = mg = \qty{78}{\kg} \cdot \qty{9.82}{\m\per\second\squared}  = \qty{765,96}{\N}
  .\end{equation*}
  We now take a global force equilibrium:
  \begin{align*}
    \sum F_y &= 0 \\
    A_y + B_y - \qty{765.96}{\N} &= 0 \\
    A_y + B_y &= \qty{765.96}{\N} 
  .\end{align*}
  Now, we take a moment equilibrium about $A$ to get:
  \begin{align*}
    \sum M_A &= 0 \\
    B_y \cdot \qty{1.5}{\m} - \qty{4}{\m} \cdot \qty{765.96}{\N} &= 0 \\
    B_y &= \frac{\qty{765.96}{\N} \cdot \qty{4}{\m} }{\qty{1.5}{\m} }  = \qty{2042,56}{\N}
  .\end{align*}
  Which means that
  \begin{equation*}
    A_y = \qty{765.96}{\N} - \qty{2042.56}{\N}   = \qty{-1276,6}{\N}
  .\end{equation*}
  Which means that the vertical forces at $A$, $B$, and $C$ are
  \begin{align*}
    A_y &= \qty{-1276,6}{\N} & B_y &= \qty{2042.56}{\N} & C_y &= \qty{-765.96}{\N} 
  .\end{align*}
  As the shear is constant on each span, the moment will be linear and the maximum moment will thus be at $B$:
  \begin{equation*}
    M_B = A_y \cdot \qty{1.5}{\m} = -\qty{1276.76}{\N} \cdot \qty{1.5}{\m}  = \qty{-1915,14}{\N.\m}
  .\end{equation*}
  
  To investigate the cross-section we break it into four parts, one top flange and three webs.
  The cross-sectional area of the top flange is:
  \begin{equation*}
    A_f = \qty{350}{\mm} \cdot \qty{20}{\mm} =  \qty{7e3}{\mm^2} 
  \end{equation*}
  and the cross-sectional areas of each of the webs are
  \begin{equation*}
    A_w = \qty{10}{\mm} \cdot \qty{30}{\mm} = \qty{300}{\mm^2} 
  .\end{equation*}
  Measuring from the top, the location of the centroid of the flange is:
  \begin{equation*}
    y_f = \qty{10}{\mm} 
  \end{equation*}
  and the centroidal location of the webs are:
  \begin{equation*}
    y_w = \qty{20}{\mm} + \frac{\qty{30}{\mm} }{2} = \qty{35}{\mm} 
  .\end{equation*}
  
  The neutral axis of the composite board is then found as:
  \begin{equation*}
    \overline{y} = \frac{\sum A_i y_i}{A_i} = \frac{\qty{7e3}{\mm^2} \cdot \qty{10}{\mm} + 3 \left( \qty{300}{\mm^2} \cdot \qty{35}{\mm}  \right)}{\qty{7e3}{\mm^2} + 3 \cdot \qty{300}{\mm^2} } = \qty{12.85}{\mm} 
  .\end{equation*}
  
  Now we will find the second moment of the board about the neutral axis.
  The flange has a second moment about its own axis of:
  \begin{equation*}
    I_{c,f} = \frac{bh^3}{12} = \frac{\qty{350}{\mm} \cdot \left( \qty{20}{\mm}  \right)^3}{12} = \qty{233333}{\mm^4}
  .\end{equation*}
  The distance from its centroidal axis to the neutral axis is
  \begin{equation*}
    d_f = \qty{12.85}{\mm} - \qty{10}{\mm} = \qty{2.85}{\mm} 
  .\end{equation*}
  So using the parallel-axis theorem the flange's second moment about the neutral axis is
  \begin{equation*}
    I_f = I_{c,f} + A_f d_f^2 = \qty{233333}{\mm^{4}} + \qty{7000}{\mm^2} \cdot \left( \qty{2.85}{\mm}  \right)^2 = \qty{290190}{\mm^4}
  .\end{equation*}
  Using the same method for the webs we get:
  \begin{align*}
    I_{c,w} &= \frac{bh^3}{12} =  \frac{\qty{10}{\mm} \cdot \left( \qty{30}{\mm}  \right)^3}{12}  = \qty{22500}{\mm^4} \\
    d_w &= \qty{35}{\mm} - \qty{12.85}{\mm}   = \qty{22,15}{\mm}\\
    I_w &= I_{c,w} + A_w d_w^2 = \qty{22500}{\mm^{4}} + \qty{300}{\mm^2} \cdot \left( \qty{22.15}{\mm}  \right)^2  = \qty{169686}{\mm^4}
  .\end{align*}
  Hence, the total $I$ about the neutral axis is:
  \begin{equation*}
    I = I_f + 3 I_w = \qty{290190}{\mm^{4}} + 3 \cdot \qty{169686}{\mm^{4}}  = \qty{799248}{\mm^4}
  .\end{equation*}

  The extreme fiber distance from the neutral axis is to the bottom
  \begin{equation*}
    c_{\mathrm{max}} = \left( \qty{20}{\mm} + \qty{30}{\mm}  \right)- \overline{y} = \qty{37.15}{\mm} 
  .\end{equation*}
  
  The magnitude of the maximum bending stress is then:
  \begin{equation*}
    \sigma_{\mathrm{max}} = \frac{M_{\mathrm{max}} c_{\mathrm{max}}}{I} = \frac{\qty{1915}{\N.\m} \cdot \qty{37.15}{\mm} }{\qty{799248}{\mm^{4}} } = \qty{89,0}{\MPa} 
  .\end{equation*}
  
  Which means the maximum strain is:
  \begin{equation*}
    \epsilon_{\mathrm{max}} = \frac{\sigma_{\mathrm{max}}}{E} = \frac{\qty{89,0}{\MPa}}{\qty{125}{\GPa} } = \num{7,12e-4} 
  .\end{equation*}
  Which occurs at the bottom-fiber at $B$.
\end{derivation}

\finalresult{\epsilon_{\mathrm{max}} = \num{7,12e-4}}
